\documentclass{../note}

\begin{document}

\begin{zadanie}{1}
Niech $f:[0,1]\rightarrow[0,1]$ będzie rosnącą odwracalną funkcją wklęsłą spełniającą $f(0)=0$ i $f(1)=1$. Udowodnij nierówność $f(x)f^{-1}(x)\le x^{2}$ dla $x\in[0,1]$.
\end{zadanie}
\begin{rozwiazanie}
Z definicji funkcji wypukłej
\[f(x) = f((1-x)\cdot 0 + x\cdot 1) \geq (1-x)f(0) + xf(1) = x\]
$f^{-1}$ jest rosnąca, zatem z powyższej nierówności wynika $f^{-1}(x) \leq x$. Mamy więc
\[x = f(f^{-1}(x)) = f\left(\frac{x-f^{-1}(x)}{x} \cdot 0 + \frac{f^{-1}(x)}{x}\cdot x\right) \geq\]\[ \frac{x-f^{-1}(x)}{x}f(0) + \frac{f^{-1}(x)}{x}\cdot f(x) = \frac{f^{-1}(x)f(x)}{x}\]
czyli $f^{-1}(x)f(x) \leq x^2$.
\end{rozwiazanie}

\begin{zadanie}{3}
Niech $f:\mathbb{R}\rightarrow\mathbb{R}$ będzie funkcją wypukłą. Udowodnij, że dla dowolnych $x,y,z\in\mathbb{R}$ prawdziwa jest nierówność
$$
\frac{f(x)+f(y)+f(z)}{3}+f\left(\frac{x+y+z}{3}\right)\ge\frac{2}{3}\left(f\left(\frac{x+y}{2}\right)+f\left(\frac{y+z}{2}\right)+f\left(\frac{z+x}{2}\right)\right)
$$
\end{zadanie}
\begin{rozwiazanie}
BSO załóżmy, że $x \leq y \leq z$ oraz $y - x > z - y$. Wtedy
\[x \leq \frac{x+y}{2} \leq \frac{x + z}{2} \leq \frac{x + y + z}{3} \leq y \leq \frac{y+z}{2} \leq z\]
Pierwsza nierówność wynika z pierwszego założenia BSO, druga też, $\frac{x + z}{2} \leq y$ z drugiego założenia BSO oraz $\frac{x + y + z}{3}$ jest ich średnią ważoną, więc jest pomiędzy nimi, a piąta i szósta wynikają z pierwszego założenia BSO. Z definicji funkcji wypukłej:
\begin{equation}
\frac{f(y)}{3} + \frac{f(z)}{3} = \frac23\left(\frac12f(y) + \frac12f(z)\right) \geq \frac23f(\frac{y+z}{2})
\end{equation}
\end{rozwiazanie}

\end{document}